% SOURCE_AUTHORITY: sga5_fr_workpass.tex, lines 14452--14505 (LF-normalized unit).
% SOURCE_SHA256: 791F4EFFC5E02832D5D77ED03518C8156D6F07E4C8238B03545DB93D883FBB28
A projeção $\pi_{\Spec(\mathcal{A})/S}:\Spec(\mathcal{A})^{(p)}\to\Spec(\mathcal{A})$, respectivamente $\pi_{\Proj(\mathcal{A})/S}$, identifica-se com o morfismo deduzido de $\mathcal{A}\to\mathcal{A}^{(p)}$. Por outro lado, $\Fr_{\Spec(\mathcal{A})/S}$, respectivamente $\Fr_{\Proj(\mathcal{A})/S}$, provém de um homomorfismo de $\OO_S$-álgebras
\[
\mathcal{A}^{(p)}=\mathcal{A}\otimes_{\OO_S}\OO_{(S,\fr_S)}\longrightarrow \mathcal{A}
\]
que transforma $a\otimes f$ em $a^p f$ para toda seção $a$ de $\mathcal{A}$ e toda seção $f$ de $\OO_S$. Para vê-lo, observa-se que o composto deste homomorfismo com $\mathcal{A}\to\mathcal{A}^{(p)}$, à direita, é a elevação à $p$-ésima potência em $\mathcal{A}$.

Considere-se o caso particular em que $\mathcal{A}$ é a álgebra simétrica $\Sym(\cE)$ de um $\OO_S$-módulo quase-coerente $\cE$; é claro que $\mathcal{A}^{(p)}\simeq\Sym(\cE^{(p)})$, logo
\[
\mathbb V(\cE^{(p)})\simeq\Spec(\mathcal{A}^{(p)})\simeq\Spec(\mathcal{A})^{(p)}=\mathbb V(\cE)^{(p)},
\]
\[
\mathbb P(\cE^{(p)})\simeq\Proj(\mathcal{A}^{(p)})\simeq\Proj(\mathcal{A})^{(p)}=\mathbb P(\cE)^{(p)}.
\]
Quanto a $\Fr_{\mathbb V(\cE)/S}$, respectivamente $\Fr_{\mathbb P(\cE)/S}$, fatora-se da seguinte maneira:
\[
\mathbb V(\cE)\longrightarrow \mathbb V(\Sym^p(\cE))\longrightarrow \mathbb V(\cE^{(p)}),
\]
respectivamente
\[
\mathbb P(\cE)\longrightarrow \mathbb P(\Sym^p(\cE))\longrightarrow \mathbb P(\cE^{(p)}),
\]
onde a primeira flecha é definida pelo morfismo canônico
\[
\Sym(\Sym^p(\cE))\longrightarrow \Sym(\cE)
\]
que procede de $\Sym^p(\cE)\to\Sym(\cE)$, enquanto a segunda é definida pelo morfismo de $\OO_S$-módulos
\[
\cE^{(p)}\longrightarrow \Sym^p(\cE)
\]
proveniente do morfismo $p$-linear $\cE\to\Sym^p(\cE)$ dado pela elevação à $p$-ésima potência simétrica das seções. No caso de $X=\mathbb P(\cE)$, calculemos o feixe fundamental $\OO_{X^{(p)}}(1)$:
\[
\OO_{X^{(p)}}(1)=(\Sym(\cE)^{(p)}(1))^{\sim}
=(\Sym(\cE)(1)^{(p)})^{\sim}
\simeq \OO_X(1)^{(p/S)}.
\]
Quando $\cE=\OO_S^n$, $\Sym(\cE)=\OO_S[T_1,\ldots,T_n]$ e $X=\mathbb V(\cE)=S[T_1,\ldots,T_n]$, respectivamente $X=\mathbb P(\cE)=\mathbb P_S^{n-1}$, é claro que $\cE^{(p)}\simeq\OO_S^n=\cE$, de modo que $X^{(p)}=X$. O morfismo de Frobenius $\Fr_{X/S}$ procede do endomorfismo da álgebra $\OO_S[T_1,\ldots,T_n]$ que transforma $T_i$ em $T_i^p$ $(i=1,\ldots,n)$.

Considere-se agora um pré-esquema $X$ afim de tipo finito, respectivamente projetivo, sobre $S$, definido por um ideal, respectivamente por um ideal homogêneo,
\[
\JJ\subset \OO_S[T_1,\ldots,T_n].
\]
Assim, $X=\Spec(\mathcal{A})$, respectivamente $X=\Proj(\mathcal{A})$, onde
\[
\mathcal{A}=\OO_S[T_1,\ldots,T_n]/\JJ.
\]
Vê-se que
\[
\mathcal{A}^{(p)}\simeq \OO_S[T_1,\ldots,T_n]/\JJ',
\]
onde $\JJ'$ é a imagem em $\OO_S[T_1,\ldots,T_n]$ de $\JJ^{(p)}=\JJ\otimes_{\OO_S}\OO_{(S,\fr_S)}$; o ideal $\JJ'$ deduz-se de $\JJ$ elevando os coeficientes dos polinômios à $p$-ésima potência; tem-se $X^{(p)}=\Spec(\mathcal{A}^{(p)})$, respectivamente $\Proj(\mathcal{A}^{(p)})$, e o morfismo de Frobenius
\[
\Fr_{X/S}:X\longrightarrow X^{(p)}
\]
é definido pela elevação à $p$-ésima potência das coordenadas $T_i$, o que determina um morfismo $\mathcal{A}^{(p)}\to\mathcal{A}$.
